%!TEX root=../w02.tex
\section{Problem 3}
\begin{itemize}
    \item \textbf{(a)}:
    This is false. 
    \begin{solution}
        Suppose, there exists five consecutive integers whose average is equal to the second-largest of those integers. 
        Let $n$ be an arbitrary integer in $\Z$. 
        We set these five integers are 
        \begin{equation*}
            n - 2, \quad n - 1, \quad n, \quad n + 1, \quad\text{and} \quad n + 2. 
        \end{equation*}

        By the assumption, we have 
        \begin{equation*}
            n + 1 = \frac{\paren{n - 2} + \paren{n - 1} + n + \paren{n + 1} + \paren{n + 2}}{5}. 
        \end{equation*}

        By expanding the numerator, we get 
        \begin{equation*}
            n + 1 = \frac{n - 2 + n - 1 + n + n + 1 + n + 2}{5}. 
        \end{equation*}

        By adding up the numerator, we get 
        \begin{equation*}
            n + 1 = \frac{5n + 0}{5} = \frac{5n}{5}. 
        \end{equation*}

        Simplifying, we get 
        \begin{equation*}
            n + 1 = n, 
        \end{equation*}

        which is a contradiction to the fact that $n + 1 \neq n$. 

        Therefore, the original claim does not hold. 
    \end{solution}

    \newpage

    \item (b) This is true. 
    \begin{solution}
        Suppose that $k$ is an arbitrary integer, and $k^2$ is the first perfect square of the three consecutive perfect squares. 

        Therefore, the three perfect squares are 
        \begin{equation*}
            k^2, \quad \paren{k + 1}^2 \quad \paren{k + 2}^2.  
        \end{equation*}

        Take average of these three perfect squares, we have 
        \begin{equation*}
            avg = \frac{k^2 + \paren{k + 1}^2 + \paren{k + 2}^2}{3}. 
        \end{equation*}

        By expanding each term and adding them together, we get 
        \begin{equation*}
            avg = \frac{k^2 + k^2 + 2k + 1 + k^2 + 4k + 4}{3} = \frac{3k^2 + 6k + 5}{3}. 
        \end{equation*}

        Then, we can split $3k^2 + 6k + 5$ into $3k^2 + 6k + 3 + 2$, we get 
        \[
            avg = \frac{3k^2 + 6k + 3 + 2}{3} = \frac{3k^2 + 6k + 3}{3} + \frac{2}{3} = k^2 + 2k + 1 + \frac{2}{3}. 
        \]

        So there exists $c = k^2 + 2k + 1$ in $\Z$ such that 
        \begin{equation*}
            avg = c + \frac{2}{3}. 
        \end{equation*}

        Thus, we proved that the average of three consecutive perfect squares is always a non-integer. 
    \end{solution}
\end{itemize}